% !TeX root = ../教材题库_全部条目_转写工作稿.tex
% Source: XJ-B1-C04-EX-P123-Y0664
% 第4章 / 4.3 对数函数 / 教材 p.123 / PDF 第130页
% OCR未校对：公式与文字均可能有误；未完成数学TeX转写；禁止直接发给学生
% 原题及相邻原解截图（核对用）：../题目原图/XJ-B1-C04-EX-P123-Y0664-01.webp
% 本文件只在转写完成后作为题干使用；OCR 条目尚不是可靠题干。
\begin{bankitem}{4.3 例题 例题12}{第4章 / 4.3 对数函数 / 教材 p.123 / PDF 第130页}{XJ-B1-C04-EX-P123-Y0664}{OCR未校对：公式与文字均可能有误；未完成数学TeX转写；禁止直接发给学生}
\begin{ocrdraft}
[教材p.123 / PDF第130页]
\par
0〔，0]证明：函数3，一log。02一2一3）在区间（3，十亠）上递减．
\par
证明记g（x）一2一2一3．
\par
设\ensuremath{{}^\circ\mathrm C}是（3，十亠）上任意两个实数，且<则
\par
一十。一2>0，
\par
所以g(x）一一2一3在（3，十亠）上递增．
\par
因而对于（3，十。。）上任意两个实数。当<，时，
\par
有
\par
又3-log上在（0，十亠）上递减，
\par
所以
\par
logag(u)>logag(v).
\par
因此，函数3，一logg(x）一log一()2一2一3）在（3，十。。）上递减．
\par
借助计算机软件可画出函数3，一lo()2一2一3）和g(x)一2一3的图象，
\par
如图4，3一2．由图也可以看出，函数y一log上（一2x一3）在（3，十。。）上递减．
\par
g(x）：是2丫一3
\par
4
\par
3
\par
2
\par
一3一2
\par
0
\par
2
\par
4
\par
y=logg在一2丫一3）
\par
图4．3一2
\end{ocrdraft}
\end{bankitem}
